\documentclass[11pt]{article}
\usepackage[T1]{fontenc}
\usepackage[utf8]{inputenc}
\usepackage{lmodern,amsmath,amssymb,amsthm,mathtools,mathrsfs}
\usepackage[margin=25mm]{geometry}
\usepackage{booktabs,longtable,array,tabularx,enumitem}
\usepackage{xcolor,hyperref,microtype}
\hypersetup{colorlinks=true,linkcolor=blue!45!black,urlcolor=blue!45!black}
\setlength{\emergencystretch}{3em}
\setlist{nosep,leftmargin=*}
\newcommand{\Ric}{\operatorname{Ric}}
\newcommand{\Rm}{\operatorname{Rm}}
\newcommand{\tr}{\operatorname{tr}}
\newcommand{\Vol}{\operatorname{Vol}}
\newcommand{\CP}{\mathbb{CP}}
\newcommand{\dd}{\,dV_g}
\newcommand{\Rplus}{\mathbf R^+}
\newcommand{\sref}[1]{\textsf{[#1]}}
\newtheorem{proposition}{Proposition}[section]
\newtheorem{question}[proposition]{Question}
\theoremstyle{remark}
\newtheorem*{scope}{Scope}
\title{Einstein Gaps on Four-Manifolds\\[4pt]
\large Four Public Problems: Overnight Results, Proof Mechanisms,\\
Remaining Connections, and Questions for Further Research}
\author{Research synthesis from Sessions S1003A01 and S1003A02}
\date{7--8 October 2026\quad (all times: China Standard Time, UTC+8)}
\begin{document}
\maketitle

\begin{abstract}
This report connects the work performed after the evening restart on
7 October to four specified problems: Einstein nonexistence on nonminimal
surfaces of general type, optimal total curvature, smooth non-Bach-flat
global minimizers, and the negative-scalar-curvature four-dimensional
Goldberg problem. It covers A01 Steps 48--55 and A02 Steps 53--64.
The principal joint advance is a background-sensitive Burns insertion law,
including the cancellation of previously uncontrolled curvature sectors and
a logarithmic saving on general K\"ahler--Einstein backgrounds. Other work
includes conditional minimizer and bubble arguments, quantitative
Eguchi--Hanson obstruction costs, zero-defect constructions, a Weyl-channel
reformulation of the lower-bound problem, and actual-field compatibility
identities. Several tempting deductions were corrected or withdrawn.
No complete solution of one of the four requested problems is recorded in
this interval. The report identifies exactly which mathematical connections
are missing, distinguishes them from execution delays, and asks for concrete
new arguments rather than another reorganization of the research programme.
\end{abstract}

\paragraph{How to use these two files.}
Read this synthesis first. The separate, self-contained
\texttt{02\_Mathematical\_Source\_Companion.tex} retains all eighteen
mathematical working notes from this interval, a Step chronology, and the ten
delivered collaboration messages. Source keys such as \sref{A02.61} below
identify these notes. The companion preserves historical assertions and their
corrections; the present report uses the later corrected statements.
``Manuscript result'' means that a proof is written in those sources, not that
this compilation has independently reproved all preceding analysis or
established literature priority. Numerical observations and proposed routes
are identified separately.

\tableofcontents

\section{The four questions and their present positions}
\label{sec:four}

The common problem is not simply to exclude an Einstein metric. It is to
understand how well its equation can be satisfied by \emph{all} allowed
metrics, what geometry realizes the optimum, and what information actual
fields contribute beyond numerical curvature budgets.

\begin{enumerate}[label=\textbf{Q\arabic*.}]
\item \textbf{Einstein metrics on nonminimal surfaces of general type.}
Let $X$ be a minimal complex surface of general type, $A=c_1^2(X)>0$, and
$M=X\# k\overline{\CP}^{\,2}$, $k\ge1$. Determine Einstein existence or
nonexistence on the underlying smooth manifold, not merely within its
K\"ahler metrics. The remaining small-blowup range is particularly important.
\item \textbf{Optimal total curvature.} Determine
\[
 p_*(M)=\inf_{g\in\operatorname{Met}(M)}
 \int_M|\Ric_g^\circ|^2\,dV_g,
\]
or a matching asymptotic law in a specified family. Equivalently determine the
total curvature infimum with the curvature-operator norm convention below.
An improved recovery construction gives an upper bound, not the infimum.
\item \textbf{A non-Bach-flat global minimizer.} Produce a smooth closed
metric that attains this unrestricted infimum and is not Bach-flat, or obtain
an exclusion theorem answering the specified existence problem. A stationary
point, an orbifold minimizer, and a smooth global minimizer on the original
smooth manifold are different objects.
\item \textbf{Negative Goldberg in dimension four.} For a closed
almost-K\"ahler four-manifold $(M,g,J,\omega)$ with
\[
 d\omega=0,\qquad \omega=g(J\cdot,\cdot),\qquad
 \Ric_g=\lambda g,\quad\lambda<0,
\]
prove $\nabla\omega=0$, or construct a closed strictly almost-K\"ahler
counterexample. A finite jet at one point is neither alternative.
\end{enumerate}

\begin{center}\small
\begin{tabularx}{\textwidth}{@{}p{.08\textwidth}p{.16\textwidth}XX@{}}
\toprule
Problem & Initial emphasis & Actual overnight work & Missing conclusion\\
\midrule
Q1 & A01 & Supply analysis; exact-Einstein entry; joint constructions;
conditional links from A02 & A new argument deciding the remaining
nonminimal Einstein candidates\\
Q2 & A02, with A01 & General-background savings; Weyl-channel and
frame-twist formulations & A lower bound for all competitors matching an
actual recovery sequence\\
Q3 & A02 & Selected minimizing sequences; smooth and orbifold bubbles;
conditional attainment & Exclusion or control of the remaining degeneration
channels and smooth attainment on the required manifold\\
Q4 & A01, later overlapping A02 & Canonical connection, finite jets,
Darboux compatibility, almost-K\"ahler identities & A genuinely new global
constraint, or a complete counterexample\\
\bottomrule
\end{tabularx}
\end{center}

The allocation was deliberately flexible. It did not amount to four
independent full-time attacks. In particular, A01 Steps 49--52 responded to
A02's question about general-background insertion, whereas A02 Steps 62--64
moved toward field identities overlapping Q4. This overlap produced a real
construction theorem, but the final Q1 and Q3 conclusions did not receive an
equally sustained direct attack throughout the interval.

\section{Conventions and the inherited mathematical backbone}

For a closed oriented four-manifold put
\[
 P=\int|\Ric^\circ|^2,\quad S=\int s^2,\quad
 R=\int|\Ric|^2=P+\frac S4,\quad
 W_\pm=\int|W^\pm|^2,\quad T_\pm=2\chi\pm3\tau.
\]
All integrals use the same metric; Weyl norms are curvature-operator norms.
Then
\begin{align}
 4\pi^2T_\pm&=\frac S{24}+2W_\pm-\frac P2,\label{eq:cw}\\
 \int|\Rm_{\rm op}|^2&=8\pi^2\chi+P.\label{eq:full}
\end{align}
Thus Q2 is a global curvature-variation problem with a topological additive
constant, not a problem of subtracting separately minimized Ricci and scalar
energies.

When $R>0$, distinguish
\[
 \mathcal J(g)=\frac{4P}{R},\qquad
 \mathcal E(g)=4-\frac{(\int s)^2}{\Vol(g)R}
 =\mathcal J(g)+\frac{S-(\int s)^2/\Vol(g)}R.
\]
We use $\alpha_{\rm rel}=\inf\mathcal E$ for the relative problem and $p_*$
for $\inf P$. Some A02 source notes call $p_*$ ``$\alpha$''; their zero-gap
and local $P$ estimates must not be inserted into formulas for
$\alpha_{\rm rel}$ without the established comparison. In particular,
$P=0$ includes Ricci-flat metrics, whereas $\mathcal J$ is undefined when
$R=0$.

In the actual monopole setting, the budget is
$p=(a_g^+)^2\ge A>0$. Earlier work, not new overnight work, supplies:
\begin{itemize}
\item the N1 critical separation, including the negative-class-square cost,
\[
 \mathcal J(g)+\left(\frac{T_+}{A}-\frac23\right)_+
 \ge \frac{c_*}{1+(-a^2)_+/(2\pi^2A)};
\]
\item same-solution coercivity of I40/I41 type,
\[
 \Sigma\le C\{\Pi+\sqrt{\Pi Y}\},\quad
 \Sigma=S-32\pi^2p,\quad \Pi=P+D_2,\quad Y=(T_+-a^2)_+;
\]
\item finite-density endpoint control, for every fixed $B<\infty$,
\[
 \chi/A\le B\quad\Longrightarrow\quad
 \kappa_\beta\ge c(B)\frac{(\beta-2)^2}{\beta^2},\qquad\beta\ge2,
\]
with its specified actual conformal supply and normalization;
\item the linear scale in the original balanced family, actual strict
sub-four recovery constructions, and the web contribution giving an analytic
negative logarithmic term inside any fixed complex-hyperbolic collar;
\item the conditional-on-FD2 two-orientation surface-bundle argument and its
Kodaira-slope consequence. This is an inherited separate application, not a
new overnight solution of Q1--Q4.
\end{itemize}
Each retains its original hypotheses. Neither the existing backbone nor its
publication is logically contingent on solving all four questions.

\subsection{What Q1 already inherits}
For the ordinary blowup family $T_+=A-k$ and $a^2=T_+$, N1 gives
\[
 \mathcal J(g)+\left(\frac13-\frac kA\right)_+\ge\varepsilon_0.
\]
Consequently $k/A>1/3-\varepsilon_0$ excludes nonzero Einstein metrics;
the positive scalar-energy budget excludes Ricci-flat ones independently.
This strengthens the classical one-third obstruction in the matching
actual-supply range. It does not settle all $k\ge1$.

The surface-bundle/Kodaira argument cannot simply be copied onto these
blowups. The usual reversed-orientation connected-sum invariants often
vanish, so that argument's second actual supplier has not been provided.
Vanishing of a familiar invariant is not itself a theorem that no actual
class exists. The exact missing input is an appropriate supplier or a
different Einstein-specific argument. \sref{A01.48}

\section{Joint advance: a general-background Burns insertion law}
\label{sec:burns}

\subsection{From a special cell to a tensorial mechanism}
The earlier analytic complex-hyperbolic construction produced
\[
 \frac{P(g_\varepsilon)}{8\pi^2}
 =1-9\varepsilon^2\log(1/\varepsilon)+O_R(\varepsilon^2).
\]
The overnight question was whether uncontrolled $W^-$ destroys this saving
on a general negative K\"ahler--Einstein surface. Initial low-degree
numerical correctors appeared to incur a $W^-$ cost. Enlarging the angular
corrector space removed it. The important conclusion was not convergence
of another decimal: an explicit geometric correction explained the
cancellation.

Write
\[
 \Rplus=W^++\frac{s}{12}I,\qquad
 \mathscr C(g)=(2W_g^+,s_g/\sqrt{12}).
\]
The exact curvature identity becomes
\[
 P(g)=\int|\mathscr C(g)|^2-8\pi^2T_+(M).
\]
The scalar-flat ASD Burns metric satisfies $\mathscr C(g_B)=0$.
This representation is why the logarithmic interaction can be computed
from the second jet of the background rather than an uncontrolled collection
of its higher derivatives.

\begin{proposition}[The corrected insertion statement in the overnight manuscripts]
Let $g$ be smooth and $P$-critical on a neighbourhood containing the complete
support of an insertion at $p$. Let $\zeta\in\Lambda_p^+$ be unit length and
$K$ a fixed degree-zero angular corrector. The matched smooth Burns insertion
on $X\#\overline{\CP}^{\,2}$ has
\begin{equation}
 P(g_\varepsilon)=P(g)+8\pi^2
 -8\pi^2\varepsilon\Rplus_p(\zeta,\zeta)
 +4\pi^2\Phi(\Rm_p,\zeta,K)\varepsilon^2\log(1/\varepsilon)
 +O(\varepsilon^2).\label{eq:insertion}
\end{equation}
When $\Ric_p^\circ=0$ and $\zeta$ is an eigenvector of $\Rplus_p$,
the optimized coefficient is
\begin{equation}
 \Phi_{\rm opt}=-\mu\left(\mu-\frac s8\right)-|T|^2,
 \qquad T=\left(\Rplus_p|_{\zeta^\perp}\right)^\circ,
 \quad\Rplus_p\zeta=\mu\zeta.\label{eq:phi}
\end{equation}
In particular, for an Einstein background and a kernel direction,
\begin{equation}
 P(g_\varepsilon)=P(g)+8\pi^2
 -4\pi^2|T|^2\varepsilon^2\log(1/\varepsilon)+O(\varepsilon^2).
 \label{eq:kernel}
\end{equation}
\end{proposition}

The Einstein-background proof is in \sref{A01.52}; the more general
critical-background remainder is in \sref{A02.59}, with Part D controlling
the earlier text. The independently obtained Ricci-sector corrector below
extends the \emph{algebraic} cancellation. A full generalized rigidity
statement must still use the correct critical-background argument; Einstein
Bianchi identities cannot be transferred merely by changing the name of
the background.

\subsection{Why the logarithm and its remainder are controlled}
The source proofs use the following sequence.
\begin{enumerate}
\item Write the exact Burns metric as $\delta+\varepsilon k_B$. The
blowup changes $T_+$ by $-1$, accounting for the $8\pi^2$ term.
\item On a rescaled annulus put $z=\varepsilon/r^2$ and
\[
 G_{z,r}=B_z+r^2(q(r,\theta)+zK(\theta)).
\]
Since $\mathscr C(B_z)=0$, the first nonzero squared-curvature density is
$r^4|C_2(z,\theta)|^2$. Third and fourth background jets enter only at
higher powers of $r$ and do not generate the logarithmic coefficient.
\item The critical annular term is a multiple of
$\varepsilon^2\,dr/r$. Integrating from a fixed multiple of
$\sqrt\varepsilon$ to a fixed outer radius produces
$\tfrac12\varepsilon^2\log(1/\varepsilon)$.
The factorization bounds the remaining terms by $O(\varepsilon^2)$.
\item The core and outer matching region contribute $O(\varepsilon^2)$.
The order-$\varepsilon$ term is a distributional boundary flux; criticality
removes the additional bulk first variation. Its corrected value is
$-8\pi^2\Rplus_p(\zeta,\zeta)$.
\end{enumerate}
Fixed radii, fixed insertion data and a fixed smooth background are essential
to this stated remainder. It is not automatically uniform on a degenerating
family or when the number of insertions diverges.

\subsection{The two cancellation formulas}
For the pure $W^-$ second jet $h_W$, A01 found the degree-zero corrector
\[
 K_W=\frac32\operatorname{sym}(h_Wk_B)
       -\frac14\tr(h_Wk_B)\delta.
\]
It satisfies
\[
 D^2(W^+,s)_\delta[h_W,k_B]+D(W^+,s)_\delta[K_W]=0.
\]
The tensor identity, together with representation-theoretic cancellation of
the integrated cross terms, explains the vanishing $W^-$ cost.
\sref{A01.50, A01.52, A02.59D}

For a trace-free symmetric tensor $E$, put
\[
 u=|x|^2,\quad q=x^TEx,\quad
 E_H=\frac{E+J^TEJ}{2},\quad q_H=x^TE_Hx,
\]
\[
 f=-q/4,\qquad V=-uEx/12+qx/6,\qquad
 K_E=\mathcal L_Vk_B-\frac q2k_B-\frac{q_H}{4u}\delta.
\]
The pure-Ricci jet is $h_E=2f\delta+\mathcal L_V\delta$.
Conformal covariance of $W^+$, diffeomorphism covariance, and first-order
Burns harmonicity of $f_\varepsilon=f-\varepsilon q_H/(8u)$ imply
\[
 D^2(W^+,s)_\delta[h_E,k_B]+D(W^+,s)_\delta[K_E]=0.
\]
Thus the cancellation is structural, not a fitted list of coefficients.
Exact evaluations at rational points provide additional computational checks
of the derived identity, not its sole proof. \sref{A01.53R}

\subsection{The exact coefficient and its limits}
In the $(s,W^+)$ sector, symmetry reduces the degree-zero corrector calculation
to ten invariant tensors. A01 recovered the rational quadratic form by exact
polarization; A02 independently calculated the same expression. In terms of
the three eigenvalues of $\Rplus$,
\[
 \Phi_{\rm opt}
 =-\frac12\mu_1(\mu_1-\mu_2-\mu_3)
  -\frac12(\mu_2-\mu_3)^2.
\]
Here $\mu_1$ corresponds to the insertion direction and
$\mu_1+\mu_2+\mu_3=s/4$.

Two corrections matter. First, the coefficient is nonpositive on the
\emph{kernel} stratum, not on the whole cone $\Rplus\le0$:
for $W^+=0$ it is $s^2/288>0$, beneath a positive first-order cost.
Second, non-eigenvector correctors include a logarithmic frame-rotation null
direction, so the eigenvector formula is not an unrestricted quadratic
minimum in every curvature regime. These corrections preserve the
K\"ahler--Einstein application. \sref{A01.52, A02.59D}

\subsection{The global upper bound that changed during the night}
For a K\"ahler--Einstein surface,
$\Rplus=\operatorname{diag}(s/4,0,0)$ in the K\"ahler frame. A unit
$\zeta\perp\omega$ gives $\mu=0$ and $|T|^2=s^2/32$. Therefore, at fixed
distinct insertion points,
\begin{equation}
 \frac{P(\widetilde g_\varepsilon)}{8\pi^2}
 =k-\frac{ks^2}{64}\varepsilon^2\log(1/\varepsilon)
   +O(\varepsilon^2)<k.\label{eq:generalke}
\end{equation}
At $s=-24$ the coefficient is $9$, recovering the special complex-hyperbolic
calculation. For $s\ne0$, the negative logarithm dominates the remainder.

The manuscript extends the competitor to a minimal general-type surface
whose canonical model has ADE points. Choose insertion balls in the regular
locus where the desingularized background equals the KE metric exactly.
Fix $\varepsilon$ so that its saving is strict; only then take the resolution
scale sufficiently small that the background $P$ is less than this saving.
This gives
\[
 p_*\bigl(X\#k\overline{\CP}^{\,2}\bigr)<8\pi^2k.
\]
In the matched fixed-complex K\"ahler comparison range of the source's
Theorem $K_\infty$, the K\"ahler infimum remains $8\pi^2k$. The new
competitors therefore lie genuinely below that restricted optimization.

This does \emph{not} decide whether $p_*>0$, whether $p_*$ is attained, or
whether $p_*$ has scale $k$, $k^2/A$, or another scale as $k/A\to0$.
Nor does a fixed-background estimate immediately give a uniform saving
on every prescribed degenerating family. \sref{A02.59, Part D.5}

\section{The variational route to Q2 and Q3}
\label{sec:variational}

\subsection{A selected minimizing sequence instead of arbitrary-sequence classification}
A02 used Anderson's perturbation of $P$, on unit volume,
\[
 Z^{2p}_\varepsilon(g)=P(g)+\varepsilon\int(1+|\Ric|^2)^p,
\]
to study a selected minimizing sequence and its generalized limit.
This is a genuine change of proof organization: one need not classify every
possible near-minimizing sequence if one can obtain a suitable selected
sequence that determines the same infimum.

It does not remove the analytic content of the selected-sequence theorem.
The consolidated record explicitly keeps separate: compact orbifold bulk
versus collapse; lower semicontinuity; bubble-scale convergence and decay;
energy equality; removable singularities; and reconstruction of the original
smooth type. The earlier drafts occasionally treated several of these as
automatic; \sref{A02.58} and \sref{A02.59D} restore their proper roles.

Under the relevant energy identity, the Euler terms telescope to
\[
 p_*(M)=\sum_aP(V_a,g_a)+\sum_jP(N_j,g_j).
\]
This equality is useful once its hypotheses hold. It is not by itself an
exclusion of orbifold absorption, partial collapse, or exotic smooth caps.

\subsection{Moving a bubble and testing its background}
The sources extend the first-order Burns response to an ASD scalar-flat AE
bubble with the required homogeneous tail. Its leading interaction is a
pairing of $\Rplus_p$ with the bubble's flux tensor, optimized over insertion
frames. Under an admissible one-attachment relocation and energy equality,
a favorable relocation would contradict minimality.

This gives conditional curvature restrictions on the bulk, such as
$\Rplus\le0$ when suitable Burns directions can be tested. The new kernel
coefficient then tests the boundary of that cone. It rules out the familiar
KE-type splitting at suitable kernel points, but it does not eliminate the
strictly negative cone. Moreover, the bulk is generally $P$-critical, not
known to be Einstein; an argument using $\delta W^+=0$ must not be applied
without that extra fact. \sref{A02.53, A02.58, A01.52, A02.59D}

\subsection{Smooth AE costs and why they do not settle all bubbles}
For scalar-flat AE bubbles, the source's topological expression includes
\[
 \frac{P(N)}{8\pi^2}\ge
 4(e-1)+4b_1+b^- -5b^+,
\]
where $e$ is the number of ends, in the stated compactification setting.
If a valid SW argument first forces $b^+=0$, this becomes a positive cost
for exceptional rank and extra ends. Burns has cost $8\pi^2$; a two-ended
conformally flat scalar-flat neck has cost $32\pi^2$.

The order matters: a nonpositive charge does not itself imply one end or
vanishing $b^+$. The one-sided SW cap argument is separately recorded as an
input rather than inferred from the standard two-sided splitting theorem.
The $k=1$ argument can then force the exceptional class into the bulk under
its smooth-cone, decay, and semicontinuity assumptions. Even there, a
homotopy-sphere cap is not automatically a standard smooth $S^4$.
\sref{A02.53--54, A02.58}

\subsection{Orbifold channels: an actual obstruction to a simple cost sum}
For an ALE end with group $\Gamma$ the two oriented charges are
\[
 c_\pm(N)=\frac2{|\Gamma|}-2\chi(N)\mp3\tau(N)\mp3\eta(\partial N).
\]
The bookkeeping uses
\[
 T_+^{\rm orb}(V)=T_+(M)+c_+(N),
\]
and nonnegative costs must be added as
$\sum_j\max\{c_+(N_j),c_-(N_j),0\}$, not as the maximum of a signed total.

Several calculations identify genuine surviving channels:
\begin{itemize}
\item The LeBrun scalar-flat metric on $\mathcal O(-n)$ has
$P=8\pi^2(n-2)^2/n$. Thus ordinary AE costs cannot simply be assigned to
arbitrary ALE ends.
\item Wahl rational-ball fillings have charge pair
$(0,-4(1-1/n^2))$ in the source convention; Ricci-flat ALE fillings can have
zero $P$.
\item The lattice and linking-form calculations exclude some proposed
absorptions, but a compatible lattice is not a realized smooth filling.
The $[5,2]$ construction establishes a stated homeomorphism, not the missing
diffeomorphism.
\item The relevant fundamental-group conclusion is a quotient by the normal
closure of boundary loops, not automatically simple connectivity.
\end{itemize}
The earlier finite-list and collapse claims were correspondingly weakened.
The remaining inequality concerns the actual non-K\"ahler bulk together
with its bubble Weyl energy; it is not supplied by the topological charges
alone. \sref{A02.55--56, A02.58}

\subsection{The precise route to a non-Bach-flat minimizer}
In four dimensions, a smooth closed $P$-critical metric has constant scalar
curvature by its conformal Euler--Lagrange equation. If it is also Bach-flat,
the remaining equation implies $s\Ric^\circ=0$. Hence in the negative-scalar
range a smooth minimizer is either Einstein or non-Bach-flat.

This is a useful bridge, not an existence proof. To answer Q3 on a candidate
$M$ one still needs smooth global attainment on $M$, exclusion of the
Einstein alternative, and the negative-scalar input. A conditional orbifold
minimum or a lower-energy competitor does not provide the first requirement.
The remaining assumptions in \sref{A02.58} are therefore mathematical
work, not merely presentation details.

\section{Zero defect, approximation, and exact existence}

\subsection{Three constructive mechanisms}
The corrected \sref{A02.59D} organizes several ways to obtain $p_*=0$.
\begin{enumerate}
\item \textbf{Einstein orbifold plus Ricci-flat ALE replacements.}
For compatible cone-end identifications and derivative-controlled ALE decay
of order $\tau>0$, the matching annulus satisfies
\[
 P_{\rm ann}\le C\bigl(t^4+(\lambda/t)^{2\tau}\bigr).
\]
Taking $t=\lambda^{\tau/(\tau+2)}$ makes the error vanish. Insertions in a
finite tree are made at each parent's native scale. This proves a
sufficient zero-defect construction, not unrestricted Einstein existence.
\item \textbf{Bounded-curvature collapse.} With the spectral bound
$-3g\le\Ric\le3g$, one has $|\Ric^\circ|^2\le36$. Thus
\[
 p_*(M)\le36\Vol_{|r|}(M)\le36\operatorname{MinVol}(M).
\]
Appropriate polarized positive-rank F-structures or explicit collapsing
sequences therefore give $p_*=0$. This is an elementary combination of
classical mechanisms, not a new general existence theorem.
\item \textbf{Partial collapse adjoining an Einstein cusp region.}
The source constructs matched doubles and fillings with vanishing neck
costs. For a real-hyperbolic torus cusp neck
\[
 g=dt^2+e^{-2f(t)}g_{\rm flat},\qquad
 |\Ric^\circ|^2=3(f'')^2,
\]
so the turning-neck cost is $O(e^{-3T})$. The corresponding
complex-hyperbolic Heisenberg calculation gives $O(e^{-4T})$; the specified
torus fillings have $P\le C_\gamma e^{-2T}$.
\end{enumerate}

The reverse description in the source is conditional on Anderson's selected
limit theorem and Einstein removable-singularity regularity. It describes a
\emph{selected} sequence with compact orbifold bulk, a complete noncompact
Einstein domain plus collapse, or total collapse. These are not mutually
exclusive topological classes, and not a classification of every zero-defect
sequence.

\subsection{What these examples do and do not settle}
Attainment of $p_*=0$ is equivalent to existence of an Einstein metric on
the smooth manifold. It is not equivalent to existence of a KE metric for a
chosen complex structure, nor to ampleness of its canonical bundle.
Likewise, an obstruction to Einstein desingularization near one prescribed
orbifold excludes that limiting regime, not all Einstein metrics on the
resolution. The source withdrew the stronger claims.

These examples clarify why Q1 and Q2 are not equivalent: nonexistence of a
precise solution need not imply a positive integral distance from it. They
also show why Q3 cannot be decided by producing arbitrarily good approximate
metrics alone.

\section{A second substantive calculation: the Eguchi--Hanson cost}
\label{sec:eh}

The corrected \sref{A02.61, Section 5} quantifies a prescribed
desingularization regime. Let an Einstein orbifold have an
$\mathbb R^4/\mathbb Z_2$ point $p$, and insert Eguchi--Hanson at scale $a$
with unit resolution direction $\zeta$. Let $o_i$ be the three trace-free
kernel tensors constructed from the $L^2$ ASD harmonic form and parallel SD
forms. In the source normalization,
\[
 \|o_i\|_{L^2(EH_a)}^2=\frac{\pi^2}{a^4},\qquad
 B_i=\int_{EH_1}\langle D\Ric[\widehat h_2],o_i\rangle
     =\pi^2(\Rplus_p\zeta)_i.
\]
The first formula is an exact radial integral; the second is a boundary flux
of the background quadratic jet. The harmonic form's displayed potential is
defined off the bolt: the form is not globally exact.

For a fixed weighted neighbourhood of the correctly extended glued metric,
with weight $0\le w<2$ and bounded size parameter $\delta$,
\begin{equation}
 P(g)\ge\pi^2a^4
       \bigl(|\Rplus_p\zeta|-\epsilon_w(a,\delta)\bigr)_+^2,
 \qquad \epsilon_w(a,\delta)\longrightarrow0.\label{eq:ehlower}
\end{equation}
The proof tests $\Ric^\circ$ against the cut-off kernel tensor, uses
Cauchy--Schwarz, and exploits that its adjoint equation kills the leading
interior error. Comparison with the \emph{exact} Einstein background cancels
the dangerous higher-jet terms. The weights make the remaining scaled
integrals tend to zero even when their unscaled integrals grow.

Using the cited Biquard--Ozuch linear solution and matching, the adapted
competitor satisfies, for globally Einstein background,
\begin{equation}
 P(g_a^*)=\pi^2a^4|\Rplus_p\zeta|^2+O(a^6).
 \label{eq:ehupper}
\end{equation}
When $\Rplus_p\zeta=0$, the stated upper remainder improves to
$O(a^8\log(1/a))$. The adapted metric belongs to a bounded neighbourhood of a
generic centre, not to the vanishing neighbourhood originally asserted.
Thus the sharp comparison requires either the adapted centre or a generic
neighbourhood large enough to contain that competitor. An additive
comparison remains meaningful when $|\Rplus_p\zeta|$ is small; a
multiplicative asymptotic requires it bounded away from zero.

\paragraph{The common invariant with Burns.}
On $\ker\Rplus$, let $u,v$ be the other eigenvalues, $u+v=\Lambda=s/4$.
Then
\[
 \Phi_{\rm opt}=-\frac{(u-v)^2}{2}
              =-\frac{\Lambda^2}{2}+2uv.
\]
The determinant $uv$ also occurs in the cited second EH obstruction.
This is a useful relationship between a cost-reducing insertion and an
Einstein desingularization obstruction. It is not a newly proved
second-order EH theorem: the relevant obstruction is already in Ozuch's
work, while an $a^8$ minimized-defect law remains unfinished.

\begin{scope}
Equations \eqref{eq:ehlower}--\eqref{eq:ehupper} concern metrics in specified
weighted neighbourhoods. They do not prove that arbitrary $P$-near-minimizers
enter those neighbourhoods, and their $a^4$ cost tends to zero. They cannot
be used as a scale-independent positive gap or to discard every orbifold
channel in Q3.
\end{scope}

\section{A global reformulation: the Weyl channel and what it does not bound}
\label{sec:weyl}

Let $M=X\#k\overline{\CP}^{\,2}$ have the actual sign-family supply used
in the source. For a conformal class use its negative Yamabe representative
$\widehat g$, normalized by $s=-4$, and write
\[
 \kappa=\frac A2-\frac3{8\pi^2}W_+,\quad
 \upsilon=\frac{Y[g]^2}{32\pi^2}-A,\quad
 p=\frac{P(\widehat g)}{8\pi^2}.
\]
Here $\kappa$ is a conformal Weyl deficit, not the endpoint parameter
$\kappa_\beta$. Equation \eqref{eq:cw} gives the exact relation
\begin{equation}
 p=k-\frac43\kappa+\frac13\upsilon,
 \qquad P(g)\ge P(\widehat g).\label{eq:channel}
\end{equation}
If $\Pi=\max(a_\varepsilon^+)^2-A\ge0$, the classical combined estimate
becomes
\begin{equation}
 2\sqrt{A+\upsilon}+\sqrt{A-2\kappa}
 \ge3\sqrt{A+\Pi}.\label{eq:lebrun}
\end{equation}
Thus, for $0\le\kappa\le A/2$,
\begin{align}
 \upsilon&\ge A\left[\left(\frac{3-\sqrt{1-2\kappa/A}}2\right)^2-1\right],\\
 p&\ge L_A(k,\kappa)
 :=k-\kappa+\frac12\left(A-\kappa-\sqrt{A^2-2A\kappa}\right).
 \label{eq:L}
\end{align}
The function decreases only up to $4A/9$ and then increases. Its minimum
recovers $k-A/3$. For $\kappa<0$, the scalar bound gives
$p\ge k-4\kappa/3$. These corrected ranges matter.

\subsection{A precise conditional quadratic-scale implication}
If a new global theorem established $\kappa[g]\le k$ for every metric,
then for $0<k/A<1/3$,
\begin{equation}
 p_*(M)\ge4\pi^2\left(A-k-\sqrt{A^2-2Ak}\right)
 =\frac{2\pi^2k^2}{A}\left(1+O(k/A)\right).\label{eq:quadratic}
\end{equation}
A stronger bound $\kappa\le(1-c)k$ would instead force a linear cost.
Neither global Weyl bound is proved. The exact reformulation identifies the
missing quantity; it is not itself the missing theorem.

\subsection{Why a proposed rigidity conclusion was withdrawn}
The original Step 62 Corollary D inferred from $p/k\to0$ that
$\kappa/k\to1$, $\upsilon-\kappa=o(k)$ and $\Pi=o(k)$. This is invalid:
\[
 p=(k-\kappa)+\frac13(\upsilon-\kappa)
\]
is not a sum of two known nonnegative terms. The formal values
$(\kappa,\upsilon)=(2k,5k)$ give $p=0$ by cancellation. They are not a
constructed metric, but they disprove that algebraic inference.

The corrected conclusion is only
\[
 \kappa\ge k-p,\qquad \upsilon=4\kappa-3k+3p,
\]
together with \eqref{eq:lebrun}. Hence small $P$ does not automatically put
all competitors near the equality case of the classical estimate.
\sref{A02.62, corrected by A02.63}

\subsection{What the existing cells reveal}
For the recorded numerical cells, $\kappa/\upsilon$ is about $0.258$--$0.265$,
only slightly above the saving threshold $1/4$, while $\kappa/k$ remains
small. Indeed
\[
 \frac pk=1-\frac\kappa k\frac{4-1/\eta}{3},\qquad
 \eta=\kappa/\upsilon.
\]
This explains their modest savings within that family. It does not show that
all geometries share its efficiency. A diffuse or differently organized
competitor could be qualitatively different. The proposed $k^2/A$ scale is
still a candidate, not a proved recovery law.

\section{Actual-field structure: Chern flux, frame twist, and zero sets}

\subsection{A Hodge--Gauss estimate}
For a smooth orthogonal almost complex structure $J$, let $|\Omega|^2=2$ and
locally write
\[
 \nabla\Omega=\alpha\otimes\phi+\beta\otimes J\phi.
\]
The canonical curvature form is
\[
 \gamma_J=\Rm_{\rm op}(\Omega)+\alpha\wedge\beta,
 \qquad[\gamma_J]=2\pi c_1(J).
\]
Its ASD part separates primitive Ricci curvature from frame twist.
Harmonic projection and the triangle inequality yield
\begin{equation}
 P(g)\ge2\left(2\pi|c_1(J)^-|-
       \|(\alpha\wedge\beta)^-\|_2\right)_+^2.
 \label{eq:twist}
\end{equation}
For a maximizing sign class with square $A-k$ and projection square
$A+\Pi$, $|c_1(J)^-|^2=k+\Pi$. Thus any sub-standard-cost metric must carry
the missing exceptional Chern flux in the twist term. This describes a
necessary mechanism, not a bound on its maximum possible size.
\sref{A02.63}

An estimate controlling that twist by $P$ for \emph{arbitrary} $J$ is false:
even an Einstein product can carry a nonparallel orthogonal $J$ in a class
whose ASD Chern projection is nonzero. A useful strengthening must specify
an actually controlled choice of field.

\subsection{The exact almost-K\"ahler balance}
For $d\Omega=0$, set $u=|\nabla\Omega|^2/4$, and use the usual
$W^+=W_1^++W_2^++W_3^+$ decomposition with transverse star-Ricci form
$\rho_*''$. The source combines the Chern-form identity and
\eqref{eq:cw} to give
\begin{equation}
 P=2\|\gamma_J^-\|_2^2+\int u(5u+s)
     +2\|\rho_*''\|_2^2+4\|W_3^+\|_2^2.\label{eq:ak}
\end{equation}
For Einstein metrics this gives
\[
 \int u(|s|-6u)=2\|\rho_*''\|_2^2+4\|W_3^+\|_2^2,
 \qquad \frac12\int u^2\ge4\pi^2|c_1(J)^-|^2.
\]
These are integral constraints, not a pointwise bound on $u$. Algebraically
optimizing them alone returns the classical obstruction, not a stronger
global lower bound. Constant-torsion equality models can be excluded by
Chern-class and period arguments, but a quantitative separation of general
degenerating competitors remains missing.

\subsection{The new obstruction to a simple SW normalization argument}
The Step 64 analysis separates a bounded, amplitude-weighted SW field from
a globally nonvanishing unit almost-complex form. Near a transverse isolated
spinor zero, the normalized field has the Hopf-map model. In its standard
Euclidean normalization,
\[
 |\nabla n|^2=8/r^2,\qquad |(\alpha\wedge\beta)^-|^2=8/r^4,
\]
so its unweighted squared-twist integral diverges logarithmically. Its
boundary Hopf class also prevents replacing it by a smooth unit field inside
the ball while leaving the exterior unchanged. The corresponding
amplitude-weighted quantities can remain integrable and tend to zero with
the core radius.

Consequently one cannot use an estimate for the weighted SW field to control
the unweighted twist in \eqref{eq:twist} without retaining zero-set boundary
terms or providing a different field. The hand calculation is written in
\sref{A02.64}; its staged local checking script had not run at the recorded
Step boundary. The general zero-set geometry is not assumed transverse.

\subsection{Why diffuse modes remain relevant}
The AAD transport identity controls variation of the torsion amplitude but
annihilates its constant mode. It does not determine the orientation or
harmonic projection of the quadratic twist. Small isolated cores can also
have vanishing first-derivative cost in the relevant weak norm.
Thus the latest note does not favor either a positive cost per exceptional
class or the proposed quadratic law. What is needed is a relation involving
the harmonic flux, actual zero/core data, and the remaining curvature
compatibility, or a genuine diffuse closed construction. Merely applying
Cauchy--Schwarz to the already available identity cannot supply it.

\section{The boundary-cost formulation}

For an exceptional region $U\cong(\#_m\overline{\CP}^{\,2})\setminus B^4$,
define $Q(U)=\int_U(4|W^+|^2+s^2/12)$ and let $\mathfrak b(\partial U)$
denote the stated Gauss--Bonnet/signature boundary contribution. Localized
Chern--Weil gives
\[
 P(U)=Q(U)+8\pi^2(m-2)+2\mathfrak b(\partial U).
\]
If it replaces a reference ball with boundary contribution matched to error
$\delta$, then
\[
 P_g(U)\ge8\pi^2m-Q_h(B)-2\delta.
\]
For a KE reference of scalar curvature $-4$, $Q_h(B)=4\Vol_h(B)$.
This applies to thick interiors as well as small model caps. It identifies
two ways the apparent fixed local cost can be avoided: replacing substantial
background volume or changing the boundary transgression.

The missing step is not another evaluation of this identity. It is a
quantitative bound for those terms for the regions supplied by actual global
competitors, or a construction exploiting their freedom. Arbitrary
near-minimizers have not been shown to contain the required matched regions.
\sref{A02.63, corrected final section}

\section{Goldberg: what the local investigation actually established}
\label{sec:goldberg}

\subsection{The transverse component remains the important unknown}
For the Einstein almost-K\"ahler setting, write $q=|\nabla\omega|^2$ and
\[
 \eta=W^+(\omega)-\left(\frac s6+\frac q4\right)\omega,
 \qquad \langle\eta,\omega\rangle=0.
\]
With the complementary trace-free block $V$, the inherited exact balance is
\[
 \int\left(|\eta|^2+2|V|^2+\frac{3q^2}{16}\right)
 =-\frac s8\int q.
\]
For $s<0$ it permits nonzero torsion. The canonical-connection curvature gives
the closed form
\[
 \beta=iF_{\mathcal A}-\frac s4\omega,
 \quad \beta^+=\eta+\frac q8\omega,
 \quad \beta^-=-\frac q8\widehat\omega.
\]
Its cup square is another expression of the same balance, not an independent
positive term. The Hodge estimate bounds the ASD Chern projection by torsion
energy, but does not force torsion to vanish. \sref{A01.48, A01.53G}

The source also records restricted observations near or on KE backgrounds.
For example, on a compact negative KE surface a constant-length SD harmonic
form in the complex orientation is the K\"ahler form up to sign: the
holomorphic $(2,0)$ part would otherwise have constant nonzero norm,
incompatible with the negative Ricci form. Uses of Einstein rigidity or
deformation theory beyond that observation retain their cited hypotheses.
These are not the general Goldberg conclusion.

\subsection{Finite jets: useful tests, not a local existence theorem}
A01's degree-four numerical engine solved the truncated compatibility,
closedness and Einstein equations for KE, Ricci-flat and negative-Einstein
test data, with residual squares of order $10^{-27}$ in the recorded runs.
The negative tests had nonzero transverse $\eta$ and generic curvature blocks.
Thus that finite order did not force the desired integrability.

Higher-degree jobs were partly blocked and later exceeded the available
memory. Their incomplete runs are not evidence for or against existence.
Even a successful finite jet would still require formal compatibility at all
orders, convergence or a local existence theorem, and ultimately the global
closed-manifold condition.

\subsection{A decisive correction to the counting argument}
An earlier formal-dimension count predicted a negative number of degrees of
freedom at degree $15$. Step 55 identifies the omitted universal relation.
For Darboux-compatible variations $hJ+Jh=0$,
\[
 (\sigma_{\Ric}(\xi)h)(J\xi,J\xi)
 -\frac12|\xi|^2\tr(\sigma_{\Ric}(\xi)h)=0.
\]
Including its prolongations changes the count to
\[
 f_{\rm corr}(N)=N^2+8N-5,\qquad N\ge4.
\]
The former negative count therefore supplied no negative-scalar obstruction.
Its coordinate-free Ricci projection becomes a differential consequence of
$\nabla\Ric=0$ after imposing Einstein, not a new equation separating the
negative case from the controls. This does not prove local existence at all
orders either. \sref{A01.55}

The remaining route is the complete actual almost-K\"ahler
Ricci/Bianchi system with transverse curvature and torsion retained. The
long Weyl--Bochner formula in Step 48 was explicitly a hand-derived candidate
requiring further checking; it is not used here as an established new global
estimate. Repeating the degree ladder without incorporating the corrected
symbol relations would repeat an invalid inference.

\section{Chronology: every Step in the requested interval}
\label{sec:chronology}

Times below are Step start times. The last Step in each Session was still
running in the captured record. A completed Step is an execution boundary,
not a completed public problem.

\begin{longtable}{@{}p{.13\textwidth}p{.13\textwidth}p{.67\textwidth}@{}}
\toprule A01 Step & Start & Mathematical trajectory\\\midrule\endhead
48 & Oct 7, 21:53 & Read the four-problem and collected-research files;
separate Q1 supply from surface-bundle supply; identify the canonical and
Weyl--Bochner Goldberg mechanisms; reply on the two-ended quotient cell.\\
49 & 22:46 & Respond to A02's general-KE question. Build an independent
background-jet evaluator; observe that the apparent $W^-$ penalty changes
when the corrector space is enlarged.\\
50 & Oct 8, 00:49 & Derive the chiral-square interpretation and explicit
$W^-$ corrector; exact cancellation; communicate the general-KE mechanism.\\
51 & 01:48 & Independent spectral checks and exact invariant reductions.
Abandon an invalid quadratic-jet calculation in the raw Ricci representation.
The proposed division of labour was not answered within this Step.\\
52 & 03:00 & Exact rational coefficient, fixed-configuration remainder,
kernel-insertion theorem; correct the whole-cone sign claim; send message
833. Distinguish Einstein from merely $P$-critical bulk.\\
53 & 04:27 & Derive the Ricci corrector from covariance and harmonicity.
Return to Goldberg finite jets and the closed canonical form. Start expensive
higher-degree calculations; prepare a collaboration message.\\
54 & 06:20 & Shell and message sending blocked. Read actual-field sources
and work on symbol/AAD compatibility. The proposed degree-15 obstruction
was subsequently corrected.\\
55 & 08:17 & Correct the Darboux count and Ricci compatibility interpretation.
Runner interruption and same-Step recovery at 09:32; mathematical message
834 sent at 09:37.\\
\bottomrule
\end{longtable}

\begin{longtable}{@{}p{.13\textwidth}p{.13\textwidth}p{.67\textwidth}@{}}
\toprule A02 Step & Order & Mathematical trajectory\\\midrule\endhead
53 & Evening & Read the new sources; specialize the selected-minimizer
entrance to $P$; develop bubble relocation and conditional attainment;
ask A01 about arbitrary $W^-$ in the KE insertion.\\
54 & Evening & Work on $k=1$ ball quotients, smooth bubble charges,
topology and the Einstein/non-Bach-flat alternative. Add corrections to
the order of the SW and topological arguments.\\
55 & Late evening & Compute orbifold charges, linking forms and T-chain
lattice conditions; attempt to exclude orbifold absorption and collapse.\\
56 & Before midnight & Correct fundamental-group and collapse deductions;
retain conditional orbifold attainment; prove a zero-energy upper
construction using compatible Ricci-flat ALE fillings; study Wahl interaction.\\
57 & After midnight & General-background numerical Burns coefficient;
independent exact-background and normal-coordinate checks; identify
the kernel saving and the need for rigorous remainder control.\\
58 & After midnight & Consolidate assumptions, verify the invariant sector,
exchange messages 831--832 with A01; retain rather than conceal the remaining
attainment assumptions; formulate EH obstruction-cost work.\\
59 & Early morning & Write the critical-background remainder and exact
coefficient; general-KE competitor; organize zero-gap mechanisms. Later
Part D repairs supersede earlier excess claims.\\
60 & Early morning & Find and correct the frozen self-dual-projector error;
check the corrected density independently; compute EH flux; pursue exact
Ricci correctors. Several symbolic runs time out or fail to produce results.\\
61 & Morning & Apply the Step 59 corrections; derive the EH kernel norm,
local lower cost and matching upper construction with revised neighbourhoods.
Prepare feedback to A01, not yet sent through the blocked shell.\\
62 & Morning & Rewrite the global problem in $(\kappa,\upsilon)$;
identify a possible Weyl bound behind the quadratic scale. A claimed
necessary equality structure is later withdrawn.\\
63 & Morning & Hodge--Gauss/frame-twist estimate, AK identity, equality-locus
exclusion and localized boundary-cost formula. Correct Step 62 and arbitrary
twist coercivity; shell and collaboration remain blocked.\\
64 & Current snapshot & Analyze diffuse torsion, uncontrolled harmonic
projection, and the Hopf singularity of normalized SW spinors. Neither
positive per-class cost nor a quadratic-scale law follows.\\
\bottomrule
\end{longtable}

\section{Collaboration, execution, and research efficiency}

\subsection{Delivered collaboration and its mathematical effect}
The group record contains ten messages in the extracted interval:
\begin{center}\small
\begin{tabularx}{\textwidth}{@{}llX@{}}
\toprule Message/time & Sender & Content and use\\\midrule
825, 22:13 & A01 & Quotient-neck topology and standard baseline; symmetry
limitation for the suggested cell.\\
826--828, 22:14--22:58 & A02 & Selected-minimizer entrance, conditional
attainment, charge corrections, general-KE question.\\
829, 00:35 & A02 & Orbifold conservation, surviving channels and construction
questions.\\
830, 01:38 & A01 & Chiral reduction, explicit $W^-$ corrector, general-KE
saving mechanism.\\
831--832, 01:56--02:51 & A02 & Independent coefficient, proposed remainder
division, exact-sector comparison and correction requests.\\
833, 04:05 & A01 & Exact formula, remainder proof, corrected sign scope,
Einstein versus $P$-critical warning.\\
834, 09:37 & A01 & Geometric Ricci corrector, corrected Goldberg count,
request concerning the remaining critical-background statement.\\
\bottomrule
\end{tabularx}
\end{center}
The proof files actually use these exchanges. This was substantive
collaboration. It was nevertheless intermittent: A02's Step 61 and Step 63
feedback drafts remained unsent at the Step 63 boundary; A01's Step 53 draft
was also blocked before its essential content was delivered in message 834.
At the 09:42 read check, A02 had retrieved through message 833, not 834.

\subsection{Execution losses are real but not the whole explanation}
The operational record identifies memory throttling of the shared Session
scope and an erroneous exclusive environment lease on a long mathematical
job. Ordinary output filenames had been classified as shared-installation
changes. A02's Step 63 explicitly records that no shell scripts and no group
messages were sent during that Step. A01 also lost execution and messaging
opportunities. Subsequent higher-degree jobs exceeded the 96 GiB hard memory
budget and did not complete. The native recovery preserved research files
and resumed A01's same Step 55 at 09:32.

These failures explain lost execution and broken exchanges. They do not
prove that repairing the system supplies a missing global mathematical
inequality. The latter difficulty was already visible in the proof records.
No percentage of the night is assigned to either cause without adequate
timing attribution.

\subsection{The mathematical reason completion has not followed}
There are four distinct issues.
\begin{enumerate}
\item \textbf{The successful construction and the missing lower bound have
different quantifiers.} A local insertion plus gluing builds competitors.
It does not require, and therefore does not supply, control over every
near-minimizer. Repeating its optimization cannot alone determine $p_*$.
\item \textbf{Several conclusions depended on assumptions not yet obtained
from the minimizing sequence.} Smooth cone points, decay, energy equality,
removability and smooth topology reconstruction cannot be silently supplied
by a numerical cost or an abstract compactness citation.
\item \textbf{Some apparent constraints were redundant or false.}
The whole-cone sign assertion, the Weyl-channel equality inference, the
degree-15 count and the unrestricted field-normalization route were removed
or corrected. This was necessary mathematical work, but did not close the
four targets.
\item \textbf{Complementarity weakened after the constructive success.}
Both Sessions spent substantial time on overlapping local curvature and
almost-K\"ahler mechanisms. That work is relevant, but direct Q1
nonexistence and Q3 smooth attainment did not advance at the same rate.
The shared lower-bound frontier remained unresolved.
\end{enumerate}
The interval therefore contains genuine theorem-level construction progress,
conditional structural results and useful elimination of invalid routes,
but not a complete answer to one of Q1--Q4. It would be equally inaccurate to
call all of it trivial or to call every new identity a major breakthrough.

\section{Specific questions for the external mathematical discussion}
\label{sec:requests}

The requested help is mathematical: identify a new connection, a usable
existing theorem with matching hypotheses, a counterexample to a proposed
estimate, or a substantially different construction. The following questions
are deliberately not instructions to follow a single route.

\begin{enumerate}
\item \textbf{Q1, exact Einstein rather than all-metric coercivity.}
Can the actual monopole sign family and $\delta W^+=0$ force a restriction
excluding $X\#k\overline{\CP}^{\,2}$ below the N1 window? Which genuinely
Einstein-specific information has not been used by the classical numerical
budgets? Alternatively, does a plausible exact construction survive the
known topological and actual-field constraints?
\item \textbf{Q2, the new Weyl information.}
Can one establish or refute a bound such as $\kappa\le k$ in a stated
small-$k/A$ category, or replace it with a coupled bound sufficient for a
matching lower estimate? Existing inequalities mostly bound
$\upsilon$ \emph{below} in terms of $\kappa$; another inequality in the same
direction may add no information. The cancellation example in
Section~\ref{sec:weyl} must remain allowed unless geometry rules it out.
\item \textbf{Q2, genuinely different competitors.}
Can diffuse deformation, ASD gluing, or interacting defects achieve a scale
smaller than the cellwise linear one? A candidate should include the actual
closed metric, topology and transition costs, not only an algebraic energy
density. The existing strict saving already rules out treating the standard
unit cost as an unrestricted lower bound.
\item \textbf{Q3, the shortest valid attainment argument.}
Which of the assumptions retained in \sref{A02.58} actually follow from the
precise selected-minimizer theory, and which require new analysis? Can one
choose a concrete manifold on which all remaining channels can be decided,
without presupposing general bubble classification or smooth Poincar\'e?
An orbifold example is useful, but is not the requested smooth example.
\item \textbf{Q4, beyond the finite-jet count.}
Does an established local existence or formal-integrability theorem decide
the negative Einstein almost-K\"ahler local system with full transverse
curvature? If local flexibility persists, which closed-manifold identity or
global geometric structure could control $\eta$ and the torsion's harmonic
component? Neither the corrected symbol count nor a recycled cup-square
identity supplies this.
\item \textbf{Joint actual-field analysis.}
Can the normalized-spinor singularity be handled by a punctured Chern-form
identity retaining its boundary Hopf charges, or by another controlled
actual field? Would this yield a new global bound, rather than another
necessary condition for low-cost metrics?
\item \textbf{The EH/Burns relation.}
Can the common determinant invariant connect an actual near-minimizing
sequence to the local sharp neighbourhood theorem, or does it mainly
describe a separate desingularization regime? Please distinguish the
source's new energy calculation from the cited known EH obstruction.
\end{enumerate}

The useful end of the discussion would be a small number of precise new
lemmas or constructions, each explaining which full problem it changes and
why its hypotheses are available. The hypotheses and even candidate scales
may change when the mathematics requires it. A complete solution of a
natural restricted case is valuable; relabeling a conditional calculation as
a solution of the full public problem is not.

\appendix
\section{Source navigation and correction precedence}

All keys refer to complete text reproduced in the companion. Paths below
are relative to the corresponding Session's \texttt{workspace/notes/}.
\small
\begin{longtable}{@{}p{.15\textwidth}p{.79\textwidth}@{}}
\toprule Key & Source and role\\\midrule\endhead
A01.48 & \path{step48_problems1_4_goldberg_bochner_and_supply.md}:
initial Q1/Q4 supply analysis and candidate identities.\\
A01.49 & \path{step49_kernel_log_coefficient_general_background.md}:
early general-background calculation, later exact improvements supersede
its provisional cost interpretation.\\
A01.50 & \path{step50_wminus_completed_square_and_general_KE_construction.md}:
explicit $W^-$ correction and chiral-square mechanism.\\
A01.52 & \path{step52_kernel_insertion_theorem.md}:
exact coefficient, Einstein-background remainder and its restrictions.\\
A01.53R & \path{step53_ric0_geometric_corrector.md}:
covariant Ricci-sector correction.\\
A01.53G & \path{step53_goldberg_jets_and_canonical_form.md}:
finite jets and canonical form; counting interpretation superseded by 55.\\
A01.55 & \path{step55_darboux_ricci_compatibility.md}:
corrected Darboux symbol and connection to the actual AAD system.\\
A02.53 & \path{step53_anderson_entrance_bubble_exclusion_attainment.md}:
selected minimizer, relocation and analytic assumptions.\\
A02.54 & \path{step54_k1_ball_quotient_attainment_theorem.md}:
conditional $k=1$ argument; read its correction section.\\
A02.55 & \path{step55_k1_orbifold_gate_charges_lattices_collapse.md}:
charges and lattice calculations; read A02.56 Section 0 first.\\
A02.56 & \path{step56_orbifold_attainment_zero_energy_bound_Tpoint_interaction.md}:
corrected orbifold argument and ALE upper constructions.\\
A02.57 & \path{step57_kernel_burns_coefficient_general_background.md}:
general coefficient computation, superseded in scope by A02.59D.\\
A02.58 & \path{step58_consolidated_hypothesis_ledger.md}:
conditional dependencies; its then-open KE remainder is advanced in 59.\\
A02.59 / 59D & \path{step59_rigorous_burns_expansion_and_zero_gap_structure.md}:
general remainder, upper construction and zero-gap mechanisms;
\textbf{Part D overrides earlier parts}.\\
A02.61 & \path{step61_EH_quantitative_obstruction.md}:
EH flux and energy cost; \textbf{Section 5 overrides Sections 2--4}.\\
A02.62 & \path{step62_weyl_channel_reduction.md}:
Weyl/Yamabe coordinates; Corollary D and other statements are corrected in 63.\\
A02.63 & \path{step63_frame_twist_flux_and_AK_identity.md}:
frame-twist, AK and boundary-cost relations;
\textbf{the later correction section controls earlier statements}.\\
A02.64 & \path{step64_diffuse_twist_AAD_and_SW_zeros.md}:
mean-mode and normalized-spinor obstructions; current incomplete Step.\\
\bottomrule
\end{longtable}
\normalsize

\paragraph{Selected computational anchors.}
The proof notes identify the original scripts and receipts, including A01
\path{step50_exact_identity.py}, \path{step52_exact_su2minus_fast.py},
\path{step53_ric0_geometric_corrector.py},
\path{step53_ak_einstein_jets.py}; and A02
\path{step59_exact_log_coefficient.py},
\path{step60_EH_obstruction_pairing.py},
\path{step61_EH_cokernel_norm.py}, and
\path{step64_hopf_spinor_twist.py}.
The companion is a mathematical text collection, not a new execution of
these computations. Pending, failed and superseded scripts retain those
statuses. In particular, a staged script is not counted as a successful
calculation.

\section{Literature used by the Sessions}

This is the focused source list used in the overnight arguments, not a new
exhaustive priority search. Original theorem numbers and assumptions are
discussed in the companion.
\begin{thebibliography}{9}
\bibitem{Anderson}
M. T. Anderson, \emph{Canonical metrics on 3-manifolds and 4-manifolds},
Asian J. Math. \textbf{10} (2006), 127--163.
\href{https://arxiv.org/abs/math/0511034}{arXiv:math/0511034}.
The selected-sequence and degeneration arguments use Theorem 4.10 and its
surrounding regularity and energy discussion.
\bibitem{LeBrun}
C. LeBrun, \emph{Ricci curvature, minimal volumes, and Seiberg--Witten theory},
\href{https://arxiv.org/abs/math/0003068}{arXiv:math/0003068}.
Classical scalar/Weyl bounds, equality analysis and recovery constructions.
\bibitem{Optimal}
C. LeBrun, \emph{Curvature Functionals, Optimal Metrics, and the Differential
Topology of 4-Manifolds},
\href{https://arxiv.org/abs/math/0404251}{arXiv:math/0404251}.
The curvature-variation background for Q2 and Q3.
\bibitem{AAD}
V. Apostolov, J. Armstrong and T. Dr\u aghici,
\emph{Local models and integrability of certain almost K\"ahler
4-manifolds}, Math. Ann. \textbf{323} (2002), 633--666.
\href{https://arxiv.org/abs/math/0109098}{arXiv:math/0109098}.
Canonical connection, Ricci/Bianchi identities and restricted integrability.
\bibitem{Blair}
C. LeBrun, \emph{Einstein metrics, symplectic minimality, and
pseudo-holomorphic curves}, Ann. Global Anal. Geom. \textbf{28}
(2005), 157--177.
\href{https://arxiv.org/abs/math/0409290}{arXiv:math/0409290}.
\bibitem{Ozuch}
T. Ozuch, \emph{Higher order obstructions to the desingularization of
Einstein metrics}. The Session's source object and \sref{A02.61} specify
Section 2.1, Lemma 2.3, Theorem 0.2, and the later corrections concerning
second-order obstructions and their hypotheses.
\end{thebibliography}

\paragraph{Snapshot and continuity.}
The extraction time is 8 October 2026, 09:58:28 CST (UTC+8).
Both Sessions were running at the snapshot; this compilation neither paused
them nor changed their goals, research messages, files or model settings.
It is an increment to the earlier manuscript and 7 October handoff, not a
replacement for the month-wide research corpus. The next mathematical update
should change an estimate, an admissible class, an existence statement, or a
missing proof connection identified above.
\end{document}
